% Lista acronimelor folosite în lucrare (ordine alfabetică).
\begin{longtable}{@{}>{\bfseries}l@{\hspace{1cm}}p{11.3cm}@{}}
API        & Application Programming Interface --- interfață de programare a aplicațiilor \\
BDF        & Backward Differentiation Formula --- formulă de diferențiere regresivă (metodă numerică pentru sisteme rigide) \\
BPTT       & Backpropagation Through Time --- propagarea înapoi a erorii în timp (antrenarea rețelelor recurente) \\
CAP        & Consistency, Availability, Partition tolerance --- consistență, disponibilitate, toleranță la partiționare (teoremă) \\
CNCF       & Cloud Native Computing Foundation \\
COVID-19   & Coronavirus Disease 2019 --- boala provocată de coronavirus (2019) \\
CPU        & Central Processing Unit --- unitate centrală de procesare \\
CSRF       & Cross-Site Request Forgery --- falsificarea cererilor inter-sit \\
CSV        & Comma-Separated Values --- valori separate prin virgulă \\
DB         & Database --- bază de date \\
DNS        & Domain Name System --- sistem de nume de domeniu \\
ECDC       & European Centre for Disease Prevention and Control --- Centrul European de Prevenire și Control al Bolilor \\
EDO        & Ecuație Diferențială Ordinară (echivalent ODE) \\
GIL        & Global Interpreter Lock --- blocaj global al interpretorului (Python) \\
GPU        & Graphics Processing Unit --- unitate de procesare grafică \\
HMAC       & Hash-based Message Authentication Code --- cod de autentificare a mesajelor bazat pe funcție de dispersie \\
HPA        & Horizontal Pod Autoscaler --- autoscaler orizontal de pod-uri (Kubernetes) \\
HTTP(S)    & HyperText Transfer Protocol (Secure) --- protocol de transfer hipertext (securizat) \\
IDOR       & Insecure Direct Object Reference --- referință directă nesecurizată la obiect \\
IP         & Internet Protocol --- protocol internet \\
IQR        & Interquartile Range --- interval intercuartilic \\
JIT        & Just-In-Time (compilation) --- compilare la momentul execuției \\
JSON       & JavaScript Object Notation --- notație de obiecte JavaScript \\
JSONB      & JSON Binary --- JSON în format binar (tip de date PostgreSQL/YugabyteDB) \\
JWT        & JSON Web Token --- token web în format JSON \\
K8s        & Kubernetes --- platformă de orchestrare a containerelor \\
LLVM       & Low Level Virtual Machine --- infrastructură de compilare (folosită de Julia) \\
LSTM       & Long Short-Term Memory --- rețea neuronală cu memorie pe termen lung și scurt \\
MAE        & Mean Absolute Error --- eroare medie absolută \\
MAPE       & Mean Absolute Percentage Error --- eroare medie procentuală absolută \\
nRMSE      & normalized Root Mean Square Error --- eroare pătratică medie radicală normalizată \\
ODE        & Ordinary Differential Equation --- ecuație diferențială ordinară (v.\ EDO) \\
OWID       & Our World in Data --- sursă de date epidemiologice deschise \\
RAM        & Random Access Memory --- memorie cu acces aleatoriu \\
RBAC       & Role-Based Access Control --- control al accesului bazat pe roluri \\
REST       & Representational State Transfer (stil arhitectural pentru API-uri web) \\
RF         & Replication Factor --- factor de replicare \\
RHS        & Right-Hand Side --- membrul drept (al unui sistem de ecuații diferențiale) \\
RK4 / RK45 & Runge-Kutta de ordinul 4 / Runge-Kutta-Fehlberg de ordinul 4(5) (metode numerice) \\
RMSE       & Root Mean Square Error --- eroare pătratică medie radicală \\
SEIR       & Susceptible-Exposed-Infected-Recovered --- model compartimental cu compartiment de latență (expuși) \\
SGD        & Stochastic Gradient Descent --- coborâre pe gradient stochastic \\
SHA-256    & Secure Hash Algorithm (256 de biți) --- algoritm de dispersie securizat \\
SIDARTHE   & model compartimental cu 8 stări: Susceptibili, Infectați, Diagnosticați, simptomatici nedepistați (\emph{Ailing}), recunoscuți (\emph{Recognized}), în pericol (\emph{Threatened}), vindecați (\emph{Healed}), decedați (\emph{Extinct}) \\
SIR        & Susceptible-Infected-Recovered --- Susceptibili-Infectați-Recuperați (model compartimental de bază) \\
SIRS       & Susceptible-Infected-Recovered-Susceptible --- model SIR cu pierderea imunității în timp \\
SIS        & Susceptible-Infected-Susceptible --- model fără imunitate durabilă \\
SPA        & Single Page Application --- aplicație web cu o singură pagină \\
SQL        & Structured Query Language --- limbaj de interogare structurat \\
TLS        & Transport Layer Security --- securitatea stratului de transport \\
UI         & User Interface --- interfață cu utilizatorul \\
URL        & Uniform Resource Locator --- localizator uniform de resurse \\
VM         & Virtual Machine --- mașină virtuală \\
YSQL       & Yugabyte SQL --- interfața SQL (compatibilă PostgreSQL) a YugabyteDB \\
\end{longtable}
